\section*{Capítulo \ref{ch:b1:species-interactions} --- Interacciones entre especies}

\begin{solution}{exo:b1:species-interactions:1}
\omterm{def:b1:species-interactions:types}{Competencia} $-/-$ (dos robles por la luz); \omterm{def:b1:species-interactions:types}{depredación} $+/-$ (lince,
liebre); \omterm{def:b1:species-interactions:types}{parasitismo} $+/-$ (garrapata, ciervo); \omterm{def:b1:species-interactions:types}{mutualismo} $+/+$ (abeja,
flor); \omterm{def:b1:species-interactions:types}{comensalismo} $+/0$ (helecho sobre la corteza); amensalismo $-/0$
(hierba bajo la sombra de un roble).
\end{solution}

\begin{solution}{exo:b1:species-interactions:2}
Dos especies que usan el mismo recurso limitante de la misma manera no
pueden coexistir indefinidamente en un ambiente estable. Las cinco reinitas
usan el mismo alimento pero en partes distintas del árbol y a horas
distintas: sus \omterm{def:b1:ecosystem-organization:niche}{nichos} difieren, así que no hay dos de ellas que lo hagan ``de
la misma manera''.
\end{solution}

\begin{solution}{exo:b1:species-interactions:3}
\omterm{def:b1:species-interactions:functional}{Tiempo de manipulación}: el tiempo de capturar, comer y digerir una presa,
durante el cual no se toma ninguna otra. Tasa de ataque: la superficie
explorada por unidad de tiempo por la probabilidad de captura. Tasa máxima
$1/h = 2$ presas por hora.
\end{solution}

\begin{solution}{exo:b1:species-interactions:4}
Una especie cuya retirada cambia la estructura de la \omterm{def:b1:ecosystem-organization:ecosystem}{comunidad} de forma
desproporcionada respecto de su abundancia, por lo general al controlar a un
competidor dominante. La abundancia mide la presencia, no el efecto; un
depredador raro que se come al que iba a ganar tiene un efecto que sus
efectivos no revelan.
\end{solution}

\begin{solution}{exo:b1:species-interactions:5}
\omterm{prop:b1:species-interactions:lotka}{Isoclina} 1: $N_1 + 0.5 N_2 = 500$ (cortes en 500 en $N_1$ y 1000 en
$N_2$); \omterm{prop:b1:species-interactions:lotka}{isoclina} 2: $N_2 + 0.6 N_1 = 400$ (400 en $N_2$, 667 en $N_1$).
Se cruzan con cada especie limitada más por sí misma ($\alpha = 0.5 <
K_1/K_2 = 1.25$; $\beta = 0.6 < K_2/K_1 = 0.8$): coexistencia estable
en $N_1 = 429$, $N_2 = 143$. Con $\alpha = 2$: la \omterm{prop:b1:species-interactions:lotka}{isoclina} 1 corta en
500 y 250, ambos dentro de los de la \omterm{prop:b1:species-interactions:lotka}{isoclina} 2 (667 y 400): la especie 2
gana siempre.
\end{solution}

\begin{solution}{exo:b1:species-interactions:6}
$f = aN/(1 + ahN)$: a 10, $2/(1 + 0.5) = 1.33$ por hora; a 100,
$20/(1 + 5) = 3.33$ por hora (máximo $1/h = 4$). Semisaturación a
$N = 1/(ah) = 20$ por metro cuadrado.
\end{solution}

\begin{solution}{exo:b1:species-interactions:7}
Con una respuesta de tipo III, la fracción de presas comidas sube con la
densidad a densidades bajas: las presas raras se toman proporcionalmente
menos y pueden recuperarse. Con una de tipo II la fracción comida es máxima
a baja densidad (allí el depredador nunca está saturado), de modo que una
\omterm{def:b1:populations:population}{población} pequeña de presas se hunde todavía más.
\end{solution}

\begin{solution}{exo:b1:species-interactions:8}
\emph{P.~caudatum} usaba las mismas bacterias en la misma columna de agua
que \emph{P.~aurelia}, y la más rápida y frugal \emph{P.~aurelia} llevó el
alimento por debajo de lo que aquella necesitaba. \emph{P.~bursaria} se
alimentaba en el fondo de recursos distintos: decidió el \omterm{def:b1:ecosystem-organization:niche}{nicho}, no la tasa
de crecimiento --- \emph{P.~bursaria} no crece más deprisa que
\emph{P.~caudatum}.
\end{solution}

\begin{solution}{exo:b1:species-interactions:9}
Los depredadores se encuentran el patrón más a menudo sin sufrir daño, lo
aprenden peor y lo atacan más: la protección cae tanto para el mimético como
para el modelo. La eficacia del mimético baja por tanto conforme se vuelve
común, una dependencia negativa de la frecuencia que lo mantiene raro
respecto del modelo.
\end{solution}

\begin{solution}{exo:b1:species-interactions:10}
Sin el depredador, la \omterm{prop:b1:species-interactions:lotka}{isoclina} del mejillón queda por fuera de la de todos
sus competidores: gana la roca. La \omterm{def:b1:species-interactions:types}{depredación} rebaja la $K$ efectiva del
mejillón (su \omterm{prop:b1:species-interactions:lotka}{isoclina} se retrae hacia dentro) hasta que corta a las de los
demás desde dentro: los competidores están ahora limitados más por sí mismos
que por el mejillón, y coexisten con él. El depredador actúa sobre el
competidor más fuerte y hace sitio para el resto.
\end{solution}

\begin{solution}{exo:b1:species-interactions:11}
Cuando el fosfato limita el crecimiento, el \qty{80}{\%} del fosfato vale
mucho más que el \qty{15}{\%} del azúcar: \omterm{def:b1:species-interactions:types}{mutualismo}. En un suelo rico en
fosfato la planta captaría ella misma el fosfato y el hongo le costaría el
\qty{15}{\%} a cambio de nada: \omterm{def:b1:species-interactions:types}{parasitismo}. Experimento: cultívense plantas
con y sin el hongo a una serie de niveles de fosfato y mídase la biomasa; la
asociación es mutualista donde las plantas micorrizadas crecen más, y
parásita donde crecen menos.
\end{solution}

\begin{solution}{exo:b1:species-interactions:12}
Exclusiones: parcelas valladas a lo largo de los ríos en las que los
uapitíes no pueden entrar, emparejadas con parcelas abiertas, replicadas en
varios valles con y sin manadas de lobos, medidas durante años antes y
después de la reintroducción; regístrense la altura y la cobertura de los
sauces, la densidad de uapitíes y el tiempo que pasan en cada parcela
(huellas, cámaras), la caza y el espesor de nieve. La cascada predice que
los sauces se recuperen en las parcelas abiertas solo donde haya lobos, al
mismo ritmo que dentro de las exclusiones; si los sauces se recuperan igual
con y sin lobos, o siguen a la nieve y a la caza en su lugar, la cascada
queda refutada o compartida con esos factores.
\end{solution}

\begin{solution}{pb:b1:species-interactions:1}
\textbf{1.} A: $N_A + 1.4 N_B = 200$, cortes en 200 ($N_A$) y 143
($N_B$). B: $N_B + 0.9 N_A = 130$, cortes en 130 ($N_B$) y 144
($N_A$).
\textbf{2.} Los cortes (200 frente a 144 en el eje $N_A$, 143
frente a 130 en el eje $N_B$) dejan la \omterm{prop:b1:species-interactions:lotka}{isoclina} de A enteramente fuera de la
de B: resolver $N_A = 200 - 1.4 N_B$ con $N_B = 130 - 0.9 N_A$ da
$N_A = -69$, de modo que las rectas se encuentran fuera del cuadrante
positivo. A gana desde cualquier punto de partida: dondequiera que B pueda
todavía crecer, A crece también y la sobrevive.
\textbf{3.} $T_A = 0.87$ días, $T_B = 1.26$ días: A, la de crecimiento más
rápido.
\textbf{4.} A: $N_A + 0.3 N_C = 200$ (200, 667); C: $N_C + 0.3 N_A =
100$ (100, 333). Cada especie se limita a sí misma más que a la otra
($0.3 < 200/100$ y $0.3 < 100/200 = 0.5$): coexistencia estable.
\textbf{5.} Un bosque no es constante (las estaciones y las perturbaciones
reinician la carrera antes de que se resuelva) y tiene muchos recursos y
lugares, de modo que las especies reparten \omterm{def:b1:ecosystem-organization:niche}{nichos} en vez de compartir uno.
\textbf{6.} $N_A = 200 - 0.3 N_C$, $N_C = 100 - 0.3 N_A = 100 - 60 +
0.09 N_C$, de donde $N_C = 44$ y $N_A = 187$.
\textbf{7.} Sube y se aplana hacia 10 aproximadamente: tipo II.
\textbf{8.} $(1/N, 1/f)$: $(0.2, 0.5)$, $(0.1, 0.303)$, $(0.05, 0.2)$,
$(0.025, 0.149)$, $(0.0125, 0.125)$: una recta.
\textbf{9.} Pendiente $(0.5 - 0.125)/(0.2 - 0.0125) = 2.0 = 1/a$, de modo
que $a =
\qty{0.5}{m^2/d}$; ordenada en el origen $0.5 - 2\times 0.2 = 0.1$
días $= h$.
\textbf{10.} Máximo $1/h = 10$ presas al día; semisaturación en $N =
1/(ah) = 20$ por metro cuadrado.
\textbf{11.} Fracción de manipulación $= f h = f/10$: más de la mitad cuando
$f
> 5$, es decir, a 40 y a 80 por metro cuadrado (5,0 a 20 es exactamente
la mitad).
\textbf{12.} El crecimiento de las presas $0.1 N (1 - N/200)$ iguala al
consumo $0.5N/(1 + 0.05N)$: $0.1(1 - N/200) = 0.5/(1 + 0.05 N)$, es decir,
$(1 - N/200)(1 + 0.05N) = 5$: $1 + 0.05N - N/200 - N^2/4000 = 5$, de donde
$N^2 - 180 N + 16\,000 = 0$, $N = 90 \pm \sqrt{8100 - 16\,000}$: sin \omterm{def:b1:flowering-plant-organization:organs}{raíz}
real --- un escarabajo por metro cuadrado come más de lo que las presas
pueden producir jamás (crecimiento máximo $rK/4 = 5$ al día, en $N = 100$,
donde el escarabajo come 8,3) y las lleva a la extinción. No hay equilibrio;
las presas solo persisten si los escarabajos son menos o si ellas tienen
refugios.
\textbf{13.} El \omterm{def:b1:species-interactions:functional}{tiempo de manipulación} se reduce a la mitad para la segunda
presa ($h' = 0.05$), de modo que en la saturación el escarabajo come hasta
20 de ellas al día: la ingesta total a densidad alta se duplica más o menos
si se pasa a la presa más rápida de manipular.
\textbf{14.} $H_{15} = \ln 15 = 2.71$. Ocho especies: $-(0.85\ln 0.85
+ 7\times 0.0214\ln 0.0214) = 0.138 + 0.576 = 0.71$.
\textbf{15.} En un factor $2.71/0.71 = 3.8$.
\textbf{16.} El $-$ de la estrella de mar sobre el mejillón mantenía débil
el $-$ del mejillón sobre sus competidores; retírese el primero y el
mejillón alcanza la densidad a la que su \omterm{def:b1:species-interactions:types}{competencia} excluye al resto.
\textbf{17.} La biomasa mide cuánto del flujo de energía retiene una
especie, no lo que hacen sus interacciones; el efecto de una \omterm{def:b1:species-interactions:keystone}{especie clave}
pasa por las especies que controla, cuya biomasa sí es grande.
\textbf{18.} Retirar la estrella de mar habría liberado a una especie rara
incapaz de apoderarse de la roca: la \omterm{def:b1:ecosystem-organization:ecosystem}{comunidad} dominada por mejillones
apenas habría cambiado, y la estrella de mar no habría sido una \omterm{def:b1:species-interactions:keystone}{especie
clave}.
\textbf{19.} Jaulas que excluyan a la estrella de mar de pequeños rodales,
con jaulas abiertas como control del efecto de la propia jaula: si los
mejillones se apoderan del sitio solo en las jaulas cerradas, el mecanismo
es la \omterm{def:b1:species-interactions:types}{depredación} sobre los mejillones; unas jaulas con la estrella de mar
presente pero con los mejillones retirados a mano permitirían comprobar si
importan las demás presas de la estrella.
\textbf{20.} Ganancia $\qty{40}{mg}$, coste $8\times 0.5 = \qty{4}{mg}$:
neto \qty{36}{mg} $=$ \qty{0.58}{kJ}.
\textbf{21.} \qty{40}{\micro g} de \qty{4}{mg}: el \qty{1}{\%} del
fotosintato diario de la flor compra unas tres transferencias de polen, la
reproducción de la planta.
\textbf{22.} Cada uno toma algo del otro, pero cada uno gana más de lo que
pierde: el signo es el del efecto neto, no el de la explotación.
\textbf{23.} Ladrón $+$, planta $-$ (néctar tomado, sin polinización): un
\omterm{def:b1:species-interactions:types}{parasitismo} del \omterm{def:b1:species-interactions:types}{mutualismo}; la abeja encuentra flores vaciadas y gana menos,
un $-$ indirecto.
\textbf{24.} Planta $\leftrightarrow$ abeja $+/+$; ladrón $\to$ planta $-$,
ladrón $\to$ abeja $-$; ave $\to$ ladrón $-$. Dos signos menos: el ave
tiene un $+$ indirecto sobre la planta y sobre la abeja.
\textbf{25.} $a = \qty{0.5}{m^2/d}$, $h = \qty{0.1}{d}$ (\qty{2.4}{h}),
máximo \num{10} presas al día.
\end{solution}
